\documentclass[10pt,a4paper]{amsart}
\usepackage[margin=19mm]{geometry}
\usepackage{amsmath,amssymb,mathtools}
\usepackage[scr=rsfs,cal=euler]{mathalfa}
\usepackage{xfrac}
\usepackage[unicode,hidelinks,bookmarks=false]{hyperref}
\newcommand{\ie}{i.e.\ }
\newcommand{\cf}{cf.\ }
\newcommand{\resp}{respectively\ }
\newcommand{\Subsection}{\subsection}
\providecommand{\nobreakdash}{\nobreak\hbox{-}}
\setlength{\emergencystretch}{2em}
\multlinegap0pt
\usepackage{amssymb}
\usepackage{mathtools}
\newtheorem{theorem}{Theorem}
\title{Time-dependent Robin heat equation via Markovian switching}
\author{Fausto Colantoni}
\date{Source: arXiv:2604.27901v1 (CC BY 4.0)}
\begin{document}
\maketitle

\noindent\emph{Source and licence.} Time-dependent Robin heat equation via Markovian switching. Version arXiv:2604.27901v1 (CC BY 4.0), distributed by arXiv under the Creative Commons Attribution 4.0 International licence, http://creativecommons.org/licenses/by/4.0/, as recorded in the arXiv licence metadata for this exact version and shown on the abstract page https://arxiv.org/abs/2604.27901v1. Full licence notice: \texttt{licenses/arxiv-2604.27901-CC-BY-4.0.txt}.\par\smallskip

\noindent\textit{Editorial scope.} Counted unit: the article's theorem thm:averaging (source0.tex lines 328--346). The article's complete original proof of it is delivered with it (source0.tex lines 348--449, one \texttt{proof} environment of 102 lines). No mathematics is written, repaired or re-derived: the statement and the proof are the article's own lines, in the article's own notation, and no retained line is substituted or re-flowed.  Retained uncounted as the source-local context the proof needs: lines 49--52.  Nothing was omitted from any retained span, so the package carries no omission record.  No retained line contains a citation, so the wrapper carries no bibliography and no bibliography entry.  No retained line contains a figure environment, an image-inclusion command or drawing code, so no image is delivered and the package declares no asset dependency.  Editorial formatting only, in addition: an AMS \texttt{amsart} preamble that restates the article's own declarations and macros used by the retained lines, namely the environments \texttt{theorem} on separate counters, together with the standard packages those lines need.  The retained environments print in this document as Theorem 1 in order of appearance, whereas the article prints the same items with the numbering of its own full document.  The complete native source of the article as distributed in the arXiv e-print for this exact version is bundled unchanged as \texttt{sources/199507/source0.tex}.
\begin{align}
	\label{robin:bc}
	\partial_n u + \kappa u = 0\quad\text{on }\partial D,
\end{align}

\begin{theorem}\label{thm:averaging}
	Let $\{\alpha_t\}_{t\geq 0}$ be a stationary ergodic Markov chain with finite state space $S \subset [0, \infty)$, independent of the reflected Brownian motion $X$. Let $\bar\alpha=\mathbf E[\alpha_0]$ and, for $\varepsilon>0$, set $\alpha^\varepsilon_t=\alpha_{t/\varepsilon}$. For $f\in C_b(\overline D)$, define:
	\[
	u^\varepsilon(t,x)=\mathbf E_x\Big[f(X_t)\exp\Big(-\int_0^t\alpha^\varepsilon_s\,dL_s\Big)\Big].
	\]
	Then for every $T>0$ and $x\in\overline D$, $u^\varepsilon(\cdot,x)$ converges to $u^0(\cdot,x)$ in $\mathbf{P}_\alpha$-probability, uniformly for $t \in [0,T]$:
	\[
	\lim_{\varepsilon\downarrow 0} \mathbf{P}_\alpha \left( \sup_{t \in [0,T]} |u^\varepsilon(t,x) - u^0(t,x)| > \delta \right) = 0 \quad \text{for every } \delta > 0,
	\]
	where $u^0(t,x)=\mathbf{E}_x[f(X_t) \exp(-\bar \alpha L_t)]$ is the unique mild solution of the heat equation with deterministic Robin boundary condition \eqref{robin:bc} with constant parameter $\bar\alpha$:
	\begin{equation}\label{P_alphabar}
		\begin{cases}
			\partial_t u^0(t,x) = \tfrac{1}{2}\Delta u^0(t,x), & x \in D, t \in (0,T], \\
			\partial_n u^0(t,\xi) + \bar \alpha u^0(t,\xi) = 0, & \xi \in \partial D, t \in (0,T], \\
			u^0(0,x) = f(x), & x \in \overline{D}.
		\end{cases}
	\end{equation}
	
\end{theorem}

\begin{proof}
	Fix $T>0$ and $x\in\overline D$. The proof is divided into several steps.
	
	Define 
	\[
	J_\varepsilon(t):=\int_0^t \alpha_{s/\varepsilon}\,dL_s,\qquad 
	M_t^\varepsilon:=\exp\bigl(-J_\varepsilon(t)\bigr),\qquad 
	M_t^0:=\exp(-\bar\alpha L_t).
	\]
	Since $f$ is bounded and $0\le M_t^\varepsilon \le1$, we have for every $t$
	\[
	|u^\varepsilon(t,x)-u^0(t,x)|\le \|f\|_\infty\,\mathbf E_x\bigl[|M_t^\varepsilon-M_t^0|\bigr]
	\le \|f\|_\infty\,\mathbf E_x\bigl[\sup_{s\le T}|M_s^\varepsilon-M_s^0|\bigr].
	\]
	Since \(\exp(-x)\) is Lipschitz on $[0,\infty)$ (with Lipschitz constant $1$) and \(J_\varepsilon, L\) are non-negative, we have
	\[
	|M_s^\varepsilon-M_s^0|\le |J_\varepsilon(s)-\bar\alpha L_s|.
	\]
	Thus, setting $\mathcal E_\varepsilon:=\sup_{t\in[0,T]}|J_\varepsilon(t)-\bar\alpha L_t|$,
	\[
	\sup_{t\in[0,T]}|u^\varepsilon(t,x)-u^0(t,x)|\le \|f\|_\infty\,\mathbf E_x[\mathcal E_\varepsilon].
	\]
	Consequently, if we prove that $\mathcal E_\varepsilon\to0$ in $\mathbf P_{\alpha,X}$-probability,
	then $\mathbf E_x[\mathcal E_\varepsilon]\to0$ in $\mathbf P_\alpha$-probability (because $\mathcal E_\varepsilon$ is bounded and $\mathbf E_{\alpha,X}[\mathcal E_\varepsilon]\to0$ by Vitali, and Tonelli gives $\mathbf E_\alpha[\mathbf E_x[\mathcal E_\varepsilon]]\to0$, whence $\mathbf E_x[\mathcal E_\varepsilon]\to0$ in $\mathbf P_\alpha$-probability). This will establish the theorem.
	
		Let $\tau_r:=\inf\{t>0:L_t>r\}$ be the right-continuous inverse of the local time.
	Since $L_t$ is continuous and nondecreasing, and it increases only when the process hits the boundary, its right-continuous inverse $\tau_r$ is strictly increasing and right-continuous, with jumps. 
	Hence, for $r<r'$, we have $\tau_r<\tau_{r'}$ whenever both are finite. We also note that \(\tau\) is independent of $\alpha$.
	Using the change of variable $r=L_s$,
	\[
	J_\varepsilon(t)=\int_0^{L_t}\alpha_{\tau_r/\varepsilon}\,dr,\qquad 
	\bar\alpha L_t = \int_0^{L_t}\bar\alpha\,dr.
	\]
	Define $I_\varepsilon(t):=\int_0^{L_t}\bigl(\alpha_{\tau_r/\varepsilon}-\bar\alpha\bigr)dr$,
	so that $|J_\varepsilon(t)-\bar\alpha L_t|=|I_\varepsilon(t)|$.
	
	For a fixed $t$, the random variable $I_\varepsilon(t)$ depends on $\alpha$ only.
	Because $\alpha$ is stationary and ergodic with finite state space,
	its autocovariance $C_\alpha(\theta)=\mathbf E[(\alpha_0-\bar\alpha)(\alpha_\theta-\bar\alpha)]$
	satisfies $C_\alpha(\theta)\to0$ as $\theta\to\infty$.
	Conditional on $\tau$ (i.e. on the realisation of $X$), we compute
	\[
	\mathbf E_\alpha\bigl[I_\varepsilon(t)^2\bigr]
	= \int_0^{L_t}\int_0^{L_t} C_\alpha\!\Bigl(\frac{|\tau_r-\tau_{r'}|}{\varepsilon}\Bigr)\,dr\,dr'.
	\]
	For almost every pair $(r,r')$ with $r\neq r'$, the strict increase of $\tau$ gives $|\tau_r-\tau_{r'}|>0$,
	hence $|\tau_r-\tau_{r'}|/\varepsilon\to\infty$ and $C_\alpha(\cdots)\to0$.
	The diagonal $\{r=r'\}$ has Lebesgue measure zero.
	Since $C_\alpha$ is bounded by $C_\alpha(0)=\operatorname{Var}(\alpha_0)$,
	the dominated convergence theorem yields $\mathbf E_\alpha[I_\varepsilon(t)^2]\to0$ for every fixed $t$,
	and therefore $I_\varepsilon(t)\to0$ in $\mathbf P_\alpha$-probability.
	Because this holds for every realisation of $\tau$ (i.e. for every $\omega_X$), integrating over $\tau$ gives $\mathbf E_{\alpha,X}[I_\varepsilon(t)^2]\to0$, and therefore $J_\varepsilon(t)\to\bar\alpha L_t$ in $\mathbf P_{\alpha,X}$-probability.
	
	For any partition $0=t_0<t_1<\dots<t_m=T$ and any $t\in[t_{k-1},t_k]$,
	the monotonicity implies
	\[
	J_\varepsilon(t_{k-1})-\bar\alpha L_{t_k}\le J_\varepsilon(t)-\bar\alpha L_t\le
	J_\varepsilon(t_k)-\bar\alpha L_{t_{k-1}}.
	\]
	By writing \(J_\varepsilon(t_{k-1})-\bar\alpha L_{t_k} = (J_\varepsilon(t_{k-1})-\bar\alpha L_{t_{k-1}}) - (\bar\alpha L_{t_{k}}-\bar\alpha L_{t_{k-1}}) \) and \(J_\varepsilon(t_{k})-\bar\alpha L_{t_{k-1}} = (J_\varepsilon(t_{k-1})-\bar\alpha L_{t_{k}}) + (\bar\alpha L_{t_{k}}-\bar\alpha L_{t_{k-1}}) \), we get
	\[
	|J_\varepsilon(t)-\bar\alpha L_t|\le \max_{i=k-1,k}|J_\varepsilon(t_i)-\bar\alpha L_{t_i}|
	+\bar\alpha(L_{t_k}-L_{t_{k-1}}).
	\]
	Taking the supremum over $t$ gives
	\[
	\mathcal E_\varepsilon \le \max_{k=0,\dots,m}|J_\varepsilon(t_k)-\bar\alpha L_{t_k}|
	+ \bar\alpha\max_{k=1,\dots,m}(L_{t_k}-L_{t_{k-1}}).
	\]
	
	Now $L_t$ is almost surely continuous on $[0,T]$.
	For any $\delta>0$ and $\eta>0$, choose a partition fine enough so that
	\[
	\mathbf P_X\Bigl(\bar\alpha\max_k(L_{t_k}-L_{t_{k-1}})>\frac{\delta}{2}\Bigr)<\frac{\eta}{2}.
	\]
	For this fixed partition, the finite set $\{t_k\}$ has the property that
	\[
	\max_k|J_\varepsilon(t_k)-\bar\alpha L_{t_k}|\to0\quad\text{in }\mathbf P_{\alpha,X}\text{-probability},
	\]
	because each term converges and the maximum of finitely many convergent sequences converges.
	Thus, for $\varepsilon$ sufficiently small,
	\[
	\mathbf P_{\alpha,X}\Bigl(\max_k|J_\varepsilon(t_k)-\bar\alpha L_{t_k}|>\frac{\delta}{2}\Bigr)<\frac{\eta}{2}.
	\]
	Combining the two estimates,
	\[
	\mathbf P_{\alpha,X}\bigl(\mathcal E_\varepsilon>\delta\bigr)<\eta
	\]
	for all small $\varepsilon$. Hence $\mathcal E_\varepsilon\to0$ in $\mathbf P_{\alpha,X}$-probability.
	
	Recall that $\mathcal E_\varepsilon \le (\|S\|_\infty+\bar\alpha)L_T$. For reflected Brownian motion in a bounded $C^2$ domain, $\mathbf E[L_T]<\infty$ (since it is finite for every finite time \(T\)). Hence the family $\{\mathcal E_\varepsilon\}$ is uniformly integrable. Together with $\mathcal E_\varepsilon\to0$ in $\mathbf P_{\alpha,X}$-probability, this implies $\mathbf E_{\alpha,X}[\mathcal E_\varepsilon]\to0$ (Vitali's convergence theorem).
	By Tonelli's theorem,
	$\mathbf E_{\alpha,X}[\mathcal E_\varepsilon] = \mathbf E_\alpha\bigl[\mathbf E_x[\mathcal E_\varepsilon]\bigr]$.
	The non‑negative random variable $\mathbf E_x[\mathcal E_\varepsilon]$ (which depends only on $\alpha$)
	has expectation tending to zero; therefore it converges to zero in $\mathbf P_\alpha$-probability.
	Consequently,
	\[
	\sup_{t\in[0,T]}|u^\varepsilon(t,x)-u^0(t,x)|\le \|f\|_\infty\,\mathbf E_x[\mathcal E_\varepsilon]
	\longrightarrow 0\quad\text{in }\mathbf P_\alpha\text{-probability},
	\]
	which completes the proof.
\end{proof}

\end{document}
